\section{Ordering, improved mass windows, and an exact charge gap}
\label{sec:mass-gap-round2}

The positive-mass hypothesis used above has two distinct roles. It is
automatic for density-product ordering, whereas free-sea four-field Wick
ordering produces the explicitly computable mass correction in
\eqref{eq:effective-mass}. We improve its uniform bound, explain why a
continuum estimate cannot hold uniformly on arbitrarily coarse momentum
lattices, and prove a quantitative charge gap. All statements in this
section retain the finite regulators and the charge convention $Q=N-D/2$.

\subsection{Density-product ordering and nonnegative effective mass}
\label{subsec:mass-gap-ordering}

Define the alternative, precisely specified Hamiltonian
\begin{equation}
 H^{\mathrm{dp}}_0=H_D+H_\gamma+H_I+
 \frac{\alpha}{2V}\sum_{k\ne0}v_k\rho_{-k}\rho_k.
 \label{eq:mass-gap-dp}
\end{equation}
The density is still centered at the free sea. The product of these two
centered densities is \emph{not} additionally four-field Wick ordered.
Nuclear source and nuclear pair terms are added exactly as before, with
the same nuclear-diagonal convention. Relative to \eqref{eq:model}, this
changes a definite one-body operator and a scalar; it is a different
ordering prescription, not an assertion of equivalence after a harmless
vacuum-energy shift.

\begin{corollary}[No mass window for density-product ordering]
\label{cor:mass-gap-dp}
For \eqref{eq:mass-gap-dp}, the canonical Gaussian-domination theorem,
the one-nucleus bound, and the many-nucleus heat-kernel bound hold for
every $\alpha\ge0$, $m>0$, and finite fermion and photon cutoffs.
There is no restriction of the form $\alpha\Lambda<m$.
The same volume and charge-admissibility qualifications apply to the
many-nucleus theorem.
\end{corollary}
\begin{proof}
The density-square representation is already present, with
$h_{\mathrm{eff}}=D_0$ and scalar $-\Tr(PD_0)$. Thus its mass coefficient
is exactly $m\id$. The transverse current adds only Dirac-vector
coefficients, so it does not change this mass. Every positivity,
componentwise Gaussian, photon-coordinate, and electrostatic step in
the preceding proofs therefore applies unchanged. The source costs
depend on the density-square coefficients and the nuclear convention,
which are the same here.
\end{proof}

\begin{proposition}[Closure at a nonnegative mass]
\label{prop:mass-gap-semidefinite}
The canonical Gaussian-domination conclusion continues to hold if the
Hermitian spatial mass coefficient is only $M\ge0$. This includes
photon-dependent Dirac-vector terms and the photon limit already treated.
The centered finite-CAR trace is then nonnegative, though it need not be
strictly positive.
\end{proposition}
\begin{proof}
Replace $M$ by $M+\varepsilon\id$. At every finite Trotter approximation,
the strictly positive trace converges, as $\varepsilon\downarrow0$, to
the desired nonnegative trace. Alternatively, at fixed regulators the
Hamiltonian changes by $\varepsilon\,\mathrm d\Gamma(\mathsf b)$, a bounded
operator of norm at most $\varepsilon D$. Ground energies in every charge
sector therefore converge in norm-perturbation estimates. Apply the
strictly positive-mass comparison before taking this limit. No uniform
bound in $D$ is needed for this fixed-regulator passage.
\end{proof}

This closure argument does not cover an indefinite $M$: adding a scalar
$\varepsilon\id$ cannot approximate a fixed matrix with a negative
eigenvalue by positive matrices. Likewise, the decreasing-$\mathsf b$-form
argument fails when $\langle f,Mf\rangle$ changes sign. The numerical
tests described below provide neither an indefinite-mass theorem nor an
indefinite-mass counterexample. We leave that extension unresolved.
There is no analogous obstruction for a mass of one uniformly negative
sign: replacing the Dirac matrix $\mathsf b$ by $-\mathsf b$ reduces
$M\le-m_*\id$ to the same positive-mass theorem, and a limit covers
$M\le0$. The unresolved case has both mass signs in the same spatial
operator.

\subsection{A uniform discrete convolution bound}
\label{subsec:mass-gap-discrete}

Write $h=2\pi/L$ and
\begin{equation}
 \delta M(p):=m-M_{\mathrm{eff}}(p)
 =\frac{\alpha m h^3}{4\pi^2}
 \sum_{\substack{r\in h\Z^3,\ |r|\le\Lambda\\r\ne p}}
 \frac{1}{|p-r|^2\sqrt{m^2+|r|^2}}.
 \label{eq:mass-gap-deficit}
\end{equation}

\begin{lemma}[All-volume logarithmic bound with a lattice correction]
\label{lem:mass-gap-layercake}
For every admissible momentum $p$ and every $L>0$,
\begin{equation}
 \delta M(p)\le \frac{\alpha h}{4\pi^2}
 +\frac{13\alpha m}{2\pi^2}\operatorname{arsinh}(\Lambda/m).
 \label{eq:mass-gap-layercake}
\end{equation}
Consequently a sufficient condition uniform over \emph{all} volumes is
\begin{equation}
 \alpha\left[\frac{\Lambda/m}{4\pi^2}
 +\frac{13}{2\pi^2}\operatorname{arsinh}(\Lambda/m)\right]<1.
 \label{eq:mass-gap-uniform-window}
\end{equation}
At $\alpha=1/137$ this admits $\Lambda/m<5168.21$, with the decimal
given only as an evaluation of the explicit inequality.
\end{lemma}
\begin{proof}
On $h\Z^3$, fill the omitted kernel value by defining $K(0)=h^{-2}$,
$K(k)=|k|^{-2}$ for $0<|k|\le2\Lambda$, and zero outside this ball.
When $h\le2\Lambda$, $K$ is radially nonincreasing on the lattice.
So is $f(r)=\mathbf1_{|r|\le\Lambda}/\sqrt{m^2+|r|^2}$.
For every two superlevel sets $A,B$ of these functions, $A$ and $B$ are
concentric lattice balls. Hence
\[
 \#(A\cap(p+B))\le\min(\#A,\#B)=\#(A\cap B).
\]
Layer-cake integration gives $(K*f)(p)\le(K*f)(0)$. Filling the diagonal
only increases the sum in \eqref{eq:mass-gap-deficit}. Thus that sum is
at most
\[
 \frac1{h^2m}+\sum_{0<|r|\le\Lambda}
 \frac1{|r|^2\sqrt{m^2+|r|^2}}.
\]
The lattice shell $|r/h|_\infty=j$ has $24j^2+2\le26j^2$ points.
Its contribution to the second sum is at most
$26/[h^2\sqrt{m^2+h^2j^2}]$. The decreasing-function integral estimate
\[
 h\sum_{j=1}^{\lfloor\Lambda/h\rfloor}
 (m^2+h^2j^2)^{-1/2}
 \le\int_0^\Lambda(m^2+t^2)^{-1/2}\,\mathrm dt
 =\operatorname{arsinh}(\Lambda/m)
\]
proves \eqref{eq:mass-gap-layercake}. If $h>\Lambda$, the fermion band
contains only $p=0$, so the actual deficit is zero. Otherwise $h\le\Lambda$
and replacing $h$ by $\Lambda$ gives \eqref{eq:mass-gap-uniform-window}.
This also disposes of the exceptional case $h>2\Lambda$ in the kernel
definition.
\end{proof}

The next estimate recovers the exact continuum logarithmic coefficient
when the momentum lattice resolves the mass scale. Its explicit error
constant is deliberately generous.

\begin{lemma}[Continuum coefficient and an explicit mesh error]
\label{lem:mass-gap-cubature}
If $h\le m$, then, uniformly in $p\in\Gamma_{\Lambda,L}$,
\begin{equation}
 \frac{\delta M(p)}{\alpha m}
 \le\frac1\pi\operatorname{arsinh}(\Lambda/m)+34\frac hm.
 \label{eq:mass-gap-cubature}
\end{equation}
\end{lemma}
\begin{proof}
Set $F_p(x)=|x-p|^{-2}(m^2+|x|^2)^{-1/2}$, $s=|r-p|$,
and $E=\sqrt{m^2+|r|^2}$. Around every lattice point use its cube of side
$h$. Remove centers with $|r|<2h$ or $|r-p|<2h$. Each ball contains at
most $27$ lattice points. Since the center $r=p$ was already omitted,
the discrete sum $h^3\sum F_p(r)$ over these centers is at most $54h/m$.

On every remaining cube, both $|x-p|\ge bs$ and
$\sqrt{m^2+|x|^2}\ge bE$, where $b=1-\sqrt3/4>9/16$.
The elementary derivative estimates
\[
 \|\nabla^2|x-p|^{-2}\|\le6|x-p|^{-4},\quad
 |\nabla|x-p|^{-2}|=2|x-p|^{-3},
\]
\[
 |\nabla(m^2+|x|^2)^{-1/2}|\le(m^2+|x|^2)^{-1},\quad
 \|\nabla^2(m^2+|x|^2)^{-1/2}\|
 \le2(m^2+|x|^2)^{-3/2}
\]
and the arithmetic--geometric mean inequality give
\[
 \sup_{\text{cube}}\|\nabla^2F_p\|
 \le b^{-5}\left(\frac8{s^4E}+\frac4{s^2E^3}\right).
\]
The integral of the linear Taylor term vanishes. Since the cube integral
of $|x-r|^2$ is $h^5/4$, its midpoint error is at most $h^5/8$ times this
Hessian bound.

We record every lattice estimate used to sum the errors:
\begin{align*}
 \sum_{r\ne p}h^3|r-p|^{-4}
 &\le\frac{26}{h}\sum_{j\ge1}j^{-2}
 =\frac{13\pi^2}{3h}<\frac{43}{h},\\
 \sum_rh^3(m^2+|r|^2)^{-3}&<\frac{62}{m^3}.
\end{align*}
For the second estimate put $a=m/h\ge1$. The origin contributes at most
$a^{-3}$ after factoring $h^{-3}$; shells $j\le a$ contribute at most
$26a^{-3}$, using $\sum_{j\le a}j^2\le a^3$; the remaining shells
contribute at most $26(4/3)a^{-3}$, using
$\sum_{j>a}j^{-4}\le a^{-4}+(3a^3)^{-1}\le4/(3a^3)$.
Their sum is $(185/3)m^{-3}<62m^{-3}$.
Cauchy--Schwarz therefore yields
\[
 \sum_{r\ne p}\frac{h^3}{s^2E^3}
 <\frac{52}{h^{1/2}m^{3/2}},\qquad
 \sum_{r\ne p}\frac{h^3}{s^4E}<\frac{43}{mh}.
\]
As $b^{-5}<(16/9)^5<18$ and $h/m\le1$, the total error over the remaining
cubes is at most $18(43+26)h/m=1242h/m$.

Their union lies in the ball of radius $\Lambda+\sqrt3h/2$. A continuous
layer-cake intersection argument for concentric balls, identical to the
one above, bounds the integral of $F_p$ over this ball by its value at
$p=0$:
\[
 4\pi\operatorname{arsinh}((\Lambda+\sqrt3h/2)/m)
 \le4\pi\operatorname{arsinh}(\Lambda/m)+11h/m.
\]
Thus $h^3\sum F_p(r)\le4\pi\operatorname{arsinh}(\Lambda/m)+1307h/m$.
Divide by $4\pi^2$ in \eqref{eq:mass-gap-deficit} and use
$1307/(4\pi^2)<34$.
\end{proof}

For $h\le m$, positivity follows, for example, from
\begin{equation}
 \alpha\left[\pi^{-1}\operatorname{arsinh}(\Lambda/m)+34h/m\right]<1.
 \label{eq:mass-gap-resolved-window}
\end{equation}
For a family with $L\ge L_0\ge2\pi/m$, replace $h/m$ by $2\pi/(mL_0)$.
Only in a sufficiently fine-mesh regime does this approach the continuum
scale $\Lambda/m\simeq\frac12e^{\pi/\alpha}$.
It is not a uniform-in-all-volumes conclusion, and the sign of this bare
ordering-induced coefficient alone is not a statement about a measured
renormalized electron mass.

\begin{proposition}[Coarse-box obstruction]
\label{prop:mass-gap-coarse}
No bound with a purely logarithmic cutoff dependence and a volume-uniform
bounded error can dominate \eqref{eq:mass-gap-deficit} for all
$L\ge\Lambda^{-1}$. In particular, the positive-mass criterion cannot
hold for all such boxes at arbitrarily large $\Lambda/m$ in the fixed
free-sea Wick model.
\end{proposition}
\begin{proof}
Choose $L=2\pi/\Lambda$, so $h=\Lambda$ and the fermion momenta are
$0,\pm\Lambda e_1,\pm\Lambda e_2,\pm\Lambda e_3$.
For $p=\Lambda e_1$, the $r=0$ term alone gives
\begin{equation}
 \delta M(\Lambda e_1)\ge\frac{\alpha\Lambda}{4\pi^2}.
 \label{eq:mass-gap-coarse-lower}
\end{equation}
Thus $M_{\mathrm{eff}}(p)<0$ whenever
$\Lambda/m>4\pi^2/\alpha$ (about $5408.54$ at $\alpha=1/137$).
This configuration satisfies $L\ge\Lambda^{-1}$. It refutes the claimed
uniform continuum estimate, but does not by itself refute Gaussian
domination at an indefinite mass.
\end{proof}

\subsection{Quantitative transfer matrices and the exact charge gap}
\label{subsec:mass-gap-charge}

We use $b>0$ for inverse temperature to avoid confusing it with the
Dirac matrix $\mathsf b$. Consider a one-particle transfer matrix $U$
formed from finitely many evolutions $e^{-s_jh_j}$ and scalar unitary
kicks, with $\sum_js_j=b$, where
\[
 h_j=\sum_{\ell=1}^3v_{j\ell}\otimes\mathsf a_\ell
       +M_j\otimes\mathsf b,\qquad M_j\ge m_*\id>0.
\]
The spatial coefficients may differ from interval to interval. Thus
this setting includes every truncated photon-coordinate history in the
preceding proof. Put $N=D/2$.

\begin{lemma}[Transfer-spectrum annulus]
\label{lem:mass-gap-annulus}
Exactly $N$ eigenvalues of $U$, counted with algebraic multiplicity,
satisfy $|\lambda|\le e^{-bm_*}$. The remaining $N$ are their partners
$1/\overline\lambda$, and have modulus at least $e^{bm_*}$.
No diagonalizability or normality of $U$ is assumed.
\end{lemma}
\begin{proof}
For any transfer eigenvector $u$ with $Uu=\lambda u$, transport it along
the successive intervals and kicks, writing $f$ for the resulting vector.
Set $B(t)=\langle f(t),\mathsf b f(t)\rangle$. Scalar kicks preserve $B$;
on a free interval,
\[
 B'(t)=-2\langle f(t),M_jf(t)\rangle
 \le-2m_*\|f(t)\|^2<0.
\]
Since $B(b)=|\lambda|^2B(0)$ and its total decrease is strict,
$|\lambda|<1$ implies $B(0)>0$, and $|\lambda|>1$ implies $B(0)<0$.
The case $|\lambda|=1$ is impossible. In the first case the whole
trajectory has $B>0$, so $B'\le-2m_*B$ and
$|\lambda|^2\le e^{-2bm_*}$. In the second case the whole trajectory
has $B<0$, so $B'\le2m_*B$ and
$|\lambda|^2\ge e^{2bm_*}$, taking care to reverse the inequality upon
division by $B(0)<0$.
Finally $U^\dagger\mathsf C_5U=\mathsf C_5$ pairs the spectrum, with
algebraic multiplicity, under $\lambda\mapsto1/\overline\lambda$.
The absence of unit-modulus eigenvalues gives exactly $N$ pairs.
\end{proof}

\begin{lemma}[Canonical Fourier-coefficient bound]
\label{lem:mass-gap-fourier}
Let $F(\theta)$ be the centered trace with the additional charge
kick $e^{i\theta Q}$. Define its Fourier coefficients by
\[
F_q=(2\pi)^{-1}\int_0^{2\pi}e^{-iq\theta}F(\theta)\,\mathrm d\theta .
\]
Then
\begin{equation}
 |F_q|\le\binom{2N}{N+|q|}e^{-bm_*|q|}F_0,
 \qquad |q|\le N.
 \label{eq:mass-gap-fourier}
\end{equation}
In particular the prefactor is independent of $b$ and of the history.
\end{lemma}
\begin{proof}
Let $\lambda_1,\ldots,\lambda_N$ be the eigenvalues inside the annulus.
Pairing eigenvalues in the determinant, and including the centering
factor $e^{-iN\theta}$, gives
\[
 F(\theta)=C\prod_{j=1}^N|1+e^{i\theta}\lambda_j|^2,
 \qquad C>0.
\]
Indeed each paired factor is
$(1+e^{i\theta}\lambda_j)(1+e^{i\theta}/\overline\lambda_j)
=e^{i\theta}\overline\lambda_j^{-1}|1+e^{i\theta}\lambda_j|^2$;
the remaining constant is positive because $F(0)>0$ by the centered
trace lemma. This argument uses determinant eigenvalues only and remains
valid for defective matrices.
If $e_j$ are the elementary symmetric polynomials in these $N$ numbers
and $r=e^{-bm_*}$, then
\[
 F_0=C\sum_{j=0}^N|e_j|^2\ge C,\qquad
 |e_j|\le\binom Nj r^j.
\]
For $q\ge0$ it follows that
\[
 |F_q|\le C r^q\sum_{j=0}^{N-q}\binom N{j+q}\binom Nj r^{2j}
 \le C r^q\binom{2N}{N+q}\le r^q\binom{2N}{N+q}F_0.
\]
The middle identity is Vandermonde's convolution. Negative $q$ follows
by complex conjugation.
\end{proof}

\begin{theorem}[Charge gap above the exact interacting vacuum]
\label{thm:mass-gap-charge}
Suppose the effective mass satisfies $M_{\mathrm{eff}}\ge m_*\id>0$.
For the fully interacting cutoff Hamiltonian, including its finite set
of nonzero transverse photon modes,
\begin{equation}
 E_q(0)-E_0(0)\ge m_*|q|,\qquad |q|\le D/2.
 \label{eq:mass-gap-charge}
\end{equation}
More generally, for the density-square shifts $\eta$ of the canonical
Gaussian-domination theorem,
\begin{equation}
 E_q(\eta)\ge E_0(0)+m_*|q|.
 \label{eq:mass-gap-shifted}
\end{equation}
For density-product ordering one may take $m_*=m$. For a merely
nonnegative mass, the conclusion reduces to the previous comparison.
\end{theorem}
\begin{proof}
At each finite componentwise Trotter approximation and photon-coordinate
cutoff, the Hubbard--Stratonovich measures and the Mehler loop kernels
are positive. A density-square shift supplies a scalar phase of modulus
one independent of the charge angle $\theta$. Integrate in $\theta$
first, apply \eqref{eq:mass-gap-fourier}, and then take the absolute value
of this source phase. Consequently
\begin{equation}
 Z_q(b,\eta)\le\binom{D}{D/2+|q|}
 e^{-bm_*|q|}Z_0(b,0).
 \label{eq:mass-gap-thermal}
\end{equation}
For fixed coordinate cutoff, the trace-norm Trotter limit established
above preserves this inequality. Letting $b\to\infty$ removes the
finite combinatorial prefactor and gives the energy comparison for the
truncated model. The two-sided form comparison used in the photon
extension gives convergence of every sector ground energy as the
coordinate cutoff is removed; its constants may depend on the fixed
regulators, as before. Passing to that limit proves the theorem.
\end{proof}

The nonneutral source-free sectors in \eqref{eq:mass-gap-charge} are
auxiliary sectors of the zero-mode-deleted Hamiltonian. They do not
describe an isolated nonneutral torus satisfying Gauss's law. The theorem
is a charge gap, not a gap above the vacuum inside the neutral sector:
it does not give a pair-production gap or exclude low-energy photons.
Inserting \eqref{eq:mass-gap-shifted} in the source-completion proof also
adds $m_*Z_{\rm tot}$ to the previously stated nuclear-neutral lower
bounds, with exactly the same electrostatic costs.

\begin{remark}[The finite-temperature prefactor cannot simply be omitted]
\label{rem:mass-gap-thermal}
Take four one-particle modes, $h=m\mathsf b$ and no density squares or
photons. With $r=e^{-bm}$, the centered twist trace is
$r^{-2}|1+r e^{i\theta}|^4$. Hence
\[
 \frac{Z_1}{Z_0}=\frac{2r+2r^3}{1+4r^2+r^4}.
\]
This exceeds $r$ whenever $1-2r^2-r^4>0$, for example at $r=1/100$.
Thus $Z_q\le e^{-bm|q|}Z_0$ is false even in the free model. Its failure
does not affect \eqref{eq:mass-gap-charge}: the finite ground-state
degeneracy factor disappears after division by inverse temperature.
\end{remark}

\subsection{Numerical audit and the precise unresolved question}
\label{subsec:mass-gap-numerics}

The supplied mass-sum script reproduces the quoted $p=0$ values
$0.831,0.977,1.013$ at $\Lambda/m=10,30,100$ and
$L=36\pi/\Lambda$. These values are consistent with the exact sum, but
this scaling fixes the number of momentum shells and makes $m/h$ tend
to zero. It does not approach the continuum integral as the cutoff
grows. Moreover its test $p=0.99\Lambda e_1$ is generally off the allowed
momentum lattice, and therefore is not an entry of $M_{\mathrm{eff}}$.
The coarse-box example \eqref{eq:mass-gap-coarse-lower} is an analytic
counterexample to the proposed uniform extrapolation.

The random finite-Fock scripts are diagnostics, not proofs of a charge
gap or of an indefinite-mass extension. The charge-gap theorem above
settles the positive-mass case without relying on them. An additional
deterministic-seed search through $20{,}000$ eight-mode transfer histories
with one positive and one negative spatial mass eigenvalue found no
negative centered determinant; that absence is not a theorem. Conversely,
even a negative determinant for one history would obstruct this
pointwise-positive proof, but would not alone disprove the integrated
Gaussian-domination inequality. The indefinite-mass question remains
open in this report. The script
\texttt{numerics/round2\_mass\_gap\_checks.py} also checks the elementary
thermal counterexample, the coarse-box mass defect, and the published
cubature constants; none of those numerical checks replaces its proof.
